\section{Background and Related Work}
\label{sec:background}

\paragraph{Terraform.} A Terraform configuration is organised in directories (modules), each holding
\texttt{.tf} files with blocks. A block is identified by its kind, type and name within its
directory, e.g.\ \texttt{resource.aws\_s3\_bucket.logs} or \texttt{var.region}. Blocks refer to
each other through expressions such as \texttt{var.region}, \texttt{local.tags},
\texttt{module.vpc.id} or \texttt{aws\_s3\_bucket.logs.arn}; a module call connects the root
configuration to the variables and outputs of another directory. These references form a graph in
which an edit can propagate, but co-changes also follow conventions that leave no trace in it, such
as copies of a configuration kept per environment or blocks that are always added together. IaC
research has covered tools, adoption practices and the quality of such
code~\cite{rahman2019iacmapping}; recommending the co-changes of an edit has received little
attention, GA-IaCRec~\cite{bessghaier2026iacrec} being the closest work.

\paragraph{Mining version histories.} ROSE~\cite{zimmermann2004rose} treats each commit as a
transaction and mines association rules of the form ``whoever changed $a$ also changed $b$'',
ranked by confidence and filtered by minimum support and confidence. On eight large open-source
projects, its top three suggestions contained a correct location in 64\% of the cases. History rules
need past co-changes of the same pair of entities, so they cannot help with new blocks. We use ROSE
both as a signal family and, with its original settings, as a baseline.

\paragraph{Search-based recommendation for IaC.} GA-IaCRec~\cite{bessghaier2026iacrec} recommends
the files that co-change with an edited Ansible or Puppet file. For each query, a genetic algorithm
searches an ordered list of files that maximises $\alpha \cdot \mathit{history} + \beta \cdot
\mathit{similarity}$ with fixed $\alpha = \beta = 0.5$; on 20 projects it outperformed other search algorithms and
association rule mining, and its ablation (each heuristic alone vs.\ both) showed
the two heuristics to be complementary. It follows the SBSE idea of recasting a software engineering
problem as a search problem~\cite{harman2001sbse}.

\paragraph{Positioning.} Table~\ref{tab:positioning} summarises the differences. Our genetic
algorithm does not search a list per query: it learns, once, how much to trust each signal, from
labelled commits of other repositories, and the learned weights are then applied to unseen ones.

\begin{table}[ht]
\centering
\small
\caption{Positioning with respect to the closest prior work.}
\label{tab:positioning}
\begin{tabular}{@{}llll@{}}
\toprule
 & ROSE~\cite{zimmermann2004rose} & GA-IaCRec~\cite{bessghaier2026iacrec} & \tool{} (ours) \\
\midrule
Language          & general-purpose code   & Ansible, Puppet              & Terraform \\
Granularity       & files, functions       & files                        & blocks + types to add \\
Signals           & history                & history, similarity          & history, graph, text \\
Combination       & --                     & fixed weights                & learned weights (GA) \\
Role of the GA    & --                     & searches a list per query    & learns the weights \\
Evaluation        & later commits          & last 10\% of commits         & unseen repositories \\
\bottomrule
\end{tabular}
\end{table}
